%% ---------------------------------------------------------------------------
\section{Introduction}\label{sec:intro}

The firing squad synchronization problem, posed by Myhill in 1957 and made
popular by Moore~\cite{Moore64}, asks for a finite automaton with the following
property. Copies of the automaton are placed in the cells of a line of
arbitrary length $n$. At time $0$ the leftmost cell, the \emph{general}, is in a
special state and all other cells are \emph{quiescent}. At each step every cell
updates its state from its own state and the states of its two neighbours. All
cells must then enter a \emph{firing} state at the same time, and none of them
earlier. The difficulty is that the cells are identical, that each sees only
its two neighbours, and that a cell with $k$ states cannot count up to $n$:
the same finite rule has to synchronize lines of every length.

No line can be synchronized before time $2n-2$: the order of the general needs
$n-1$ steps to reach the far end of the line, and the news that the far end
exists needs $n-1$ more steps to come back. Goto~\cite{Goto62} and
Waksman~\cite{Waksman66} showed that time $2n-2$ can be achieved, and a
solution that synchronizes every line of length $n\ge2$ at time $2n-2$ is
called a \emph{minimal-time solution}. Since then, the least number of states
of a minimal-time solution has been a benchmark question of cellular automata
theory~\cite{Mazoyer86,Umeo09,Correa17}. Waksman's solution has $16$ states,
Balzer's $8$~\cite{Balzer67}, Gerken's $7$~\cite{Gerken87} and Mazoyer's
$6$~\cite{Mazoyer87}. In the other direction, Balzer~\cite{Balzer67} reported
a computer search showing that no four-state minimal-time solution exists;
Yun\`es~\cite{Yunes93} repeated the search and observed that the lines of
lengths $2,\dots,9$ suffice, and Sanders~\cite{Sanders94} found Balzer's
search incomplete and confirmed the result with a corrected one. Many
six-state minimal-time solutions are known today~\cite{Clergue18,NguyenMaignan20};
solutions with at most five states are known for restricted families of
lengths, such as the powers of two, some of them in non-minimal
time~\cite{UmeoYanagihara07,Yunes08}, and four-state protocols synchronize the
rings whose lengths are powers of two~\cite{UmeoKamikawaYunes09}. Whether a
five-state minimal-time solution exists has been open since Mazoyer's
construction~\cite{Umeo09,Correa17}.

\paragraph{Why lower bounds are hard}
A five-state rule is a table of $96$ entries (\cref{sec:lines}), two of which
are fixed by quiescence, so there are $5^{94}>10^{65}$ rules. A rule can be
refuted by a single line that it fails to synchronize, but a minimal-time
solution has to synchronize every line, so a proof that five states do not
suffice must exclude all rules. As there are finitely many rules, it suffices
to find one length $N$ such that no five-state rule synchronizes the lines of
lengths $2,\dots,N$; a negative answer can thus consist of finitely many
conditions. The difficulty is to find conditions that are strong enough and
cheap enough to check. This paper proves such conditions for every number of
states and measures how far they reach for four and five states.

\paragraph{The idea}
Until the echo of its right end comes back, a line cannot tell that it is
finite: it follows the infinite \emph{half-line}, which is started in the same
way but has no right end (\cref{lem:agree}). \Cref{fig:line} shows this for
Mazoyer's solution. A minimal-time solution therefore consists of a single
space-time diagram of the half-line, shared by all lengths, and, for each length
$n$, a \emph{reflected triangle} next to the right end, in which the line reacts
to its border. Everything that the reflected triangle of the line of length $n$
learns from the half-line is written on two anti-diagonals of the half-line
diagram (\cref{lem:determinacy}). Conditions on the half-line alone are cheap
to check and apply to all lengths at once.

\begin{figure}[t]
\centering
%% Figure: Mazoyer's solution on the line of length 24 and on the half-line, and where they differ
%% (cells from figures.py).
\begin{tikzpicture}[x=1.28mm,y=1.28mm]
  % (a) the line of length n = 24
  \begin{scope}
    %% generated by figures.py; do not edit
\fill[stG] (0,0) rectangle (1,-1);
\fill[stA] (0,-1) rectangle (1,-2);
\fill[stC] (1,-1) rectangle (2,-2);
\fill[stG] (0,-2) rectangle (1,-3);
\fill[stB] (1,-2) rectangle (2,-3);
\fill[stA] (2,-2) rectangle (3,-3);
\fill[stG] (0,-3) rectangle (1,-4);
\fill[stC] (1,-3) rectangle (2,-4);
\fill[stG] (2,-3) rectangle (4,-4);
\fill[stG] (0,-4) rectangle (1,-5);
\fill[stB] (1,-4) rectangle (2,-5);
\fill[stA] (2,-4) rectangle (3,-5);
\fill[stB] (3,-4) rectangle (4,-5);
\fill[stC] (4,-4) rectangle (5,-5);
\fill[stG] (0,-5) rectangle (1,-6);
\fill[stC] (1,-5) rectangle (2,-6);
\fill[stG] (2,-5) rectangle (3,-6);
\fill[stC] (4,-5) rectangle (5,-6);
\fill[stA] (5,-5) rectangle (6,-6);
\fill[stG] (0,-6) rectangle (1,-7);
\fill[stB] (1,-6) rectangle (2,-7);
\fill[stA] (2,-6) rectangle (3,-7);
\fill[stA] (4,-6) rectangle (6,-7);
\fill[stG] (6,-6) rectangle (7,-7);
\fill[stG] (0,-7) rectangle (1,-8);
\fill[stC] (1,-7) rectangle (2,-8);
\fill[stG] (2,-7) rectangle (3,-8);
\fill[stA] (4,-7) rectangle (5,-8);
\fill[stB] (5,-7) rectangle (7,-8);
\fill[stC] (7,-7) rectangle (8,-8);
\fill[stG] (0,-8) rectangle (1,-9);
\fill[stB] (1,-8) rectangle (2,-9);
\fill[stA] (2,-8) rectangle (3,-9);
\fill[stB] (5,-8) rectangle (6,-9);
\fill[stC] (6,-8) rectangle (8,-9);
\fill[stA] (8,-8) rectangle (9,-9);
\fill[stG] (0,-9) rectangle (1,-10);
\fill[stC] (1,-9) rectangle (2,-10);
\fill[stG] (2,-9) rectangle (4,-10);
\fill[stC] (6,-9) rectangle (7,-10);
\fill[stA] (7,-9) rectangle (9,-10);
\fill[stG] (9,-9) rectangle (10,-10);
\fill[stG] (0,-10) rectangle (1,-11);
\fill[stB] (1,-10) rectangle (2,-11);
\fill[stA] (2,-10) rectangle (3,-11);
\fill[stB] (3,-10) rectangle (4,-11);
\fill[stC] (4,-10) rectangle (5,-11);
\fill[stA] (6,-10) rectangle (8,-11);
\fill[stB] (8,-10) rectangle (10,-11);
\fill[stC] (10,-10) rectangle (11,-11);
\fill[stG] (0,-11) rectangle (1,-12);
\fill[stC] (1,-11) rectangle (2,-12);
\fill[stG] (2,-11) rectangle (3,-12);
\fill[stC] (4,-11) rectangle (5,-12);
\fill[stA] (6,-11) rectangle (7,-12);
\fill[stB] (7,-11) rectangle (9,-12);
\fill[stC] (9,-11) rectangle (11,-12);
\fill[stA] (11,-11) rectangle (12,-12);
\fill[stG] (0,-12) rectangle (1,-13);
\fill[stB] (1,-12) rectangle (2,-13);
\fill[stA] (2,-12) rectangle (3,-13);
\fill[stC] (4,-12) rectangle (5,-13);
\fill[stB] (7,-12) rectangle (8,-13);
\fill[stC] (8,-12) rectangle (10,-13);
\fill[stA] (10,-12) rectangle (12,-13);
\fill[stG] (12,-12) rectangle (13,-13);
\fill[stG] (0,-13) rectangle (1,-14);
\fill[stC] (1,-13) rectangle (2,-14);
\fill[stG] (2,-13) rectangle (3,-14);
\fill[stC] (4,-13) rectangle (5,-14);
\fill[stA] (5,-13) rectangle (6,-14);
\fill[stC] (8,-13) rectangle (9,-14);
\fill[stA] (9,-13) rectangle (11,-14);
\fill[stB] (11,-13) rectangle (13,-14);
\fill[stC] (13,-13) rectangle (14,-14);
\fill[stG] (0,-14) rectangle (1,-15);
\fill[stB] (1,-14) rectangle (2,-15);
\fill[stA] (2,-14) rectangle (3,-15);
\fill[stA] (4,-14) rectangle (6,-15);
\fill[stG] (6,-14) rectangle (7,-15);
\fill[stA] (8,-14) rectangle (10,-15);
\fill[stB] (10,-14) rectangle (12,-15);
\fill[stC] (12,-14) rectangle (14,-15);
\fill[stA] (14,-14) rectangle (15,-15);
\fill[stG] (0,-15) rectangle (1,-16);
\fill[stC] (1,-15) rectangle (2,-16);
\fill[stG] (2,-15) rectangle (3,-16);
\fill[stA] (4,-15) rectangle (5,-16);
\fill[stB] (5,-15) rectangle (7,-16);
\fill[stA] (8,-15) rectangle (9,-16);
\fill[stB] (9,-15) rectangle (11,-16);
\fill[stC] (11,-15) rectangle (13,-16);
\fill[stA] (13,-15) rectangle (15,-16);
\fill[stG] (15,-15) rectangle (16,-16);
\fill[stG] (0,-16) rectangle (1,-17);
\fill[stB] (1,-16) rectangle (2,-17);
\fill[stA] (2,-16) rectangle (3,-17);
\fill[stB] (5,-16) rectangle (6,-17);
\fill[stG] (6,-16) rectangle (7,-17);
\fill[stB] (9,-16) rectangle (10,-17);
\fill[stC] (10,-16) rectangle (12,-17);
\fill[stA] (12,-16) rectangle (14,-17);
\fill[stB] (14,-16) rectangle (16,-17);
\fill[stC] (16,-16) rectangle (17,-17);
\fill[stG] (0,-17) rectangle (1,-18);
\fill[stC] (1,-17) rectangle (2,-18);
\fill[stG] (2,-17) rectangle (4,-18);
\fill[stB] (5,-17) rectangle (7,-18);
\fill[stC] (7,-17) rectangle (8,-18);
\fill[stC] (10,-17) rectangle (11,-18);
\fill[stA] (11,-17) rectangle (13,-18);
\fill[stB] (13,-17) rectangle (15,-18);
\fill[stC] (15,-17) rectangle (17,-18);
\fill[stA] (17,-17) rectangle (18,-18);
\fill[stG] (0,-18) rectangle (1,-19);
\fill[stB] (1,-18) rectangle (2,-19);
\fill[stA] (2,-18) rectangle (3,-19);
\fill[stB] (3,-18) rectangle (4,-19);
\fill[stB] (5,-18) rectangle (6,-19);
\fill[stC] (6,-18) rectangle (8,-19);
\fill[stA] (8,-18) rectangle (9,-19);
\fill[stA] (10,-18) rectangle (12,-19);
\fill[stB] (12,-18) rectangle (14,-19);
\fill[stC] (14,-18) rectangle (16,-19);
\fill[stA] (16,-18) rectangle (18,-19);
\fill[stG] (18,-18) rectangle (19,-19);
\fill[stG] (0,-19) rectangle (1,-20);
\fill[stC] (1,-19) rectangle (2,-20);
\fill[stG] (2,-19) rectangle (4,-20);
\fill[stC] (6,-19) rectangle (7,-20);
\fill[stA] (7,-19) rectangle (9,-20);
\fill[stA] (10,-19) rectangle (11,-20);
\fill[stB] (11,-19) rectangle (13,-20);
\fill[stC] (13,-19) rectangle (15,-20);
\fill[stA] (15,-19) rectangle (17,-20);
\fill[stB] (17,-19) rectangle (19,-20);
\fill[stC] (19,-19) rectangle (20,-20);
\fill[stG] (0,-20) rectangle (1,-21);
\fill[stB] (1,-20) rectangle (2,-21);
\fill[stA] (2,-20) rectangle (3,-21);
\fill[stB] (3,-20) rectangle (4,-21);
\fill[stC] (4,-20) rectangle (5,-21);
\fill[stA] (6,-20) rectangle (9,-21);
\fill[stB] (11,-20) rectangle (12,-21);
\fill[stC] (12,-20) rectangle (14,-21);
\fill[stA] (14,-20) rectangle (16,-21);
\fill[stB] (16,-20) rectangle (18,-21);
\fill[stC] (18,-20) rectangle (20,-21);
\fill[stA] (20,-20) rectangle (21,-21);
\fill[stG] (0,-21) rectangle (1,-22);
\fill[stC] (1,-21) rectangle (2,-22);
\fill[stG] (2,-21) rectangle (3,-22);
\fill[stC] (4,-21) rectangle (5,-22);
\fill[stA] (6,-21) rectangle (9,-22);
\fill[stG] (9,-21) rectangle (10,-22);
\fill[stC] (12,-21) rectangle (13,-22);
\fill[stA] (13,-21) rectangle (15,-22);
\fill[stB] (15,-21) rectangle (17,-22);
\fill[stC] (17,-21) rectangle (19,-22);
\fill[stA] (19,-21) rectangle (21,-22);
\fill[stG] (21,-21) rectangle (22,-22);
\fill[stG] (0,-22) rectangle (1,-23);
\fill[stB] (1,-22) rectangle (2,-23);
\fill[stA] (2,-22) rectangle (3,-23);
\fill[stC] (4,-22) rectangle (5,-23);
\fill[stA] (6,-22) rectangle (8,-23);
\fill[stB] (8,-22) rectangle (10,-23);
\fill[stC] (10,-22) rectangle (11,-23);
\fill[stA] (12,-22) rectangle (14,-23);
\fill[stB] (14,-22) rectangle (16,-23);
\fill[stC] (16,-22) rectangle (18,-23);
\fill[stA] (18,-22) rectangle (20,-23);
\fill[stB] (20,-22) rectangle (22,-23);
\fill[stC] (22,-22) rectangle (23,-23);
\fill[stG] (0,-23) rectangle (1,-24);
\fill[stC] (1,-23) rectangle (2,-24);
\fill[stG] (2,-23) rectangle (3,-24);
\fill[stC] (4,-23) rectangle (5,-24);
\fill[stA] (6,-23) rectangle (7,-24);
\fill[stB] (7,-23) rectangle (9,-24);
\fill[stC] (9,-23) rectangle (11,-24);
\fill[stA] (12,-23) rectangle (13,-24);
\fill[stB] (13,-23) rectangle (15,-24);
\fill[stC] (15,-23) rectangle (17,-24);
\fill[stA] (17,-23) rectangle (19,-24);
\fill[stB] (19,-23) rectangle (21,-24);
\fill[stC] (21,-23) rectangle (23,-24);
\fill[stG] (23,-23) rectangle (24,-24);
\fill[stG] (0,-24) rectangle (1,-25);
\fill[stB] (1,-24) rectangle (2,-25);
\fill[stA] (2,-24) rectangle (3,-25);
\fill[stC] (4,-24) rectangle (5,-25);
\fill[stB] (7,-24) rectangle (8,-25);
\fill[stC] (8,-24) rectangle (11,-25);
\fill[stB] (13,-24) rectangle (14,-25);
\fill[stC] (14,-24) rectangle (16,-25);
\fill[stA] (16,-24) rectangle (18,-25);
\fill[stB] (18,-24) rectangle (20,-25);
\fill[stC] (20,-24) rectangle (22,-25);
\fill[stB] (22,-24) rectangle (23,-25);
\fill[stA] (23,-24) rectangle (24,-25);
\fill[stG] (0,-25) rectangle (1,-26);
\fill[stC] (1,-25) rectangle (2,-26);
\fill[stG] (2,-25) rectangle (3,-26);
\fill[stC] (4,-25) rectangle (5,-26);
\fill[stA] (5,-25) rectangle (6,-26);
\fill[stC] (8,-25) rectangle (11,-26);
\fill[stA] (11,-25) rectangle (12,-26);
\fill[stC] (14,-25) rectangle (15,-26);
\fill[stA] (15,-25) rectangle (17,-26);
\fill[stB] (17,-25) rectangle (19,-26);
\fill[stC] (19,-25) rectangle (21,-26);
\fill[stB] (21,-25) rectangle (22,-26);
\fill[stA] (22,-25) rectangle (23,-26);
\fill[stC] (23,-25) rectangle (24,-26);
\fill[stG] (0,-26) rectangle (1,-27);
\fill[stB] (1,-26) rectangle (2,-27);
\fill[stA] (2,-26) rectangle (3,-27);
\fill[stA] (4,-26) rectangle (6,-27);
\fill[stG] (6,-26) rectangle (7,-27);
\fill[stC] (8,-26) rectangle (10,-27);
\fill[stA] (10,-26) rectangle (12,-27);
\fill[stG] (12,-26) rectangle (13,-27);
\fill[stA] (14,-26) rectangle (16,-27);
\fill[stB] (16,-26) rectangle (18,-27);
\fill[stC] (18,-26) rectangle (20,-27);
\fill[stB] (20,-26) rectangle (21,-27);
\fill[stA] (21,-26) rectangle (22,-27);
\fill[stC] (22,-26) rectangle (23,-27);
\fill[stB] (23,-26) rectangle (24,-27);
\fill[stG] (0,-27) rectangle (1,-28);
\fill[stC] (1,-27) rectangle (2,-28);
\fill[stG] (2,-27) rectangle (3,-28);
\fill[stA] (4,-27) rectangle (5,-28);
\fill[stB] (5,-27) rectangle (7,-28);
\fill[stC] (8,-27) rectangle (9,-28);
\fill[stA] (9,-27) rectangle (11,-28);
\fill[stB] (11,-27) rectangle (13,-28);
\fill[stA] (14,-27) rectangle (15,-28);
\fill[stB] (15,-27) rectangle (17,-28);
\fill[stC] (17,-27) rectangle (19,-28);
\fill[stB] (19,-27) rectangle (20,-28);
\fill[stA] (20,-27) rectangle (21,-28);
\fill[stC] (21,-27) rectangle (22,-28);
\fill[stB] (22,-27) rectangle (23,-28);
\fill[stG] (0,-28) rectangle (1,-29);
\fill[stB] (1,-28) rectangle (2,-29);
\fill[stA] (2,-28) rectangle (3,-29);
\fill[stB] (5,-28) rectangle (6,-29);
\fill[stG] (6,-28) rectangle (7,-29);
\fill[stA] (8,-28) rectangle (10,-29);
\fill[stB] (10,-28) rectangle (12,-29);
\fill[stG] (12,-28) rectangle (13,-29);
\fill[stB] (15,-28) rectangle (16,-29);
\fill[stC] (16,-28) rectangle (18,-29);
\fill[stB] (18,-28) rectangle (19,-29);
\fill[stA] (19,-28) rectangle (20,-29);
\fill[stC] (20,-28) rectangle (21,-29);
\fill[stB] (21,-28) rectangle (22,-29);
\fill[stG] (0,-29) rectangle (1,-30);
\fill[stC] (1,-29) rectangle (2,-30);
\fill[stG] (2,-29) rectangle (4,-30);
\fill[stB] (5,-29) rectangle (7,-30);
\fill[stA] (8,-29) rectangle (9,-30);
\fill[stB] (9,-29) rectangle (13,-30);
\fill[stC] (13,-29) rectangle (14,-30);
\fill[stC] (16,-29) rectangle (17,-30);
\fill[stB] (17,-29) rectangle (18,-30);
\fill[stA] (18,-29) rectangle (19,-30);
\fill[stC] (19,-29) rectangle (20,-30);
\fill[stB] (20,-29) rectangle (21,-30);
\fill[stG] (0,-30) rectangle (1,-31);
\fill[stB] (1,-30) rectangle (2,-31);
\fill[stA] (2,-30) rectangle (3,-31);
\fill[stB] (3,-30) rectangle (4,-31);
\fill[stB] (5,-30) rectangle (6,-31);
\fill[stG] (6,-30) rectangle (7,-31);
\fill[stB] (9,-30) rectangle (12,-31);
\fill[stC] (12,-30) rectangle (14,-31);
\fill[stA] (14,-30) rectangle (15,-31);
\fill[stG] (16,-30) rectangle (17,-31);
\fill[stA] (17,-30) rectangle (18,-31);
\fill[stC] (18,-30) rectangle (19,-31);
\fill[stB] (19,-30) rectangle (20,-31);
\fill[stG] (0,-31) rectangle (1,-32);
\fill[stC] (1,-31) rectangle (2,-32);
\fill[stG] (2,-31) rectangle (4,-32);
\fill[stB] (5,-31) rectangle (7,-32);
\fill[stC] (7,-31) rectangle (8,-32);
\fill[stB] (9,-31) rectangle (11,-32);
\fill[stC] (11,-31) rectangle (13,-32);
\fill[stA] (13,-31) rectangle (15,-32);
\fill[stC] (15,-31) rectangle (16,-32);
\fill[stG] (16,-31) rectangle (17,-32);
\fill[stC] (17,-31) rectangle (18,-32);
\fill[stB] (18,-31) rectangle (19,-32);
\fill[stG] (0,-32) rectangle (1,-33);
\fill[stB] (1,-32) rectangle (2,-33);
\fill[stA] (2,-32) rectangle (3,-33);
\fill[stB] (3,-32) rectangle (4,-33);
\fill[stB] (5,-32) rectangle (6,-33);
\fill[stC] (6,-32) rectangle (8,-33);
\fill[stB] (9,-32) rectangle (10,-33);
\fill[stC] (10,-32) rectangle (12,-33);
\fill[stA] (12,-32) rectangle (14,-33);
\fill[stC] (14,-32) rectangle (15,-33);
\fill[stB] (15,-32) rectangle (16,-33);
\fill[stG] (16,-32) rectangle (17,-33);
\fill[stB] (17,-32) rectangle (18,-33);
\fill[stG] (0,-33) rectangle (1,-34);
\fill[stC] (1,-33) rectangle (2,-34);
\fill[stG] (2,-33) rectangle (4,-34);
\fill[stC] (6,-33) rectangle (8,-34);
\fill[stC] (10,-33) rectangle (11,-34);
\fill[stA] (11,-33) rectangle (13,-34);
\fill[stC] (13,-33) rectangle (14,-34);
\fill[stB] (14,-33) rectangle (15,-34);
\fill[stG] (16,-33) rectangle (17,-34);
\fill[stC] (17,-33) rectangle (18,-34);
\fill[stG] (0,-34) rectangle (1,-35);
\fill[stB] (1,-34) rectangle (2,-35);
\fill[stA] (2,-34) rectangle (3,-35);
\fill[stB] (3,-34) rectangle (4,-35);
\fill[stC] (4,-34) rectangle (5,-35);
\fill[stC] (6,-34) rectangle (8,-35);
\fill[stA] (8,-34) rectangle (9,-35);
\fill[stA] (10,-34) rectangle (12,-35);
\fill[stC] (12,-34) rectangle (13,-35);
\fill[stB] (13,-34) rectangle (14,-35);
\fill[stG] (16,-34) rectangle (17,-35);
\fill[stB] (17,-34) rectangle (18,-35);
\fill[stA] (18,-34) rectangle (19,-35);
\fill[stG] (0,-35) rectangle (1,-36);
\fill[stC] (1,-35) rectangle (2,-36);
\fill[stG] (2,-35) rectangle (3,-36);
\fill[stC] (4,-35) rectangle (5,-36);
\fill[stC] (6,-35) rectangle (7,-36);
\fill[stA] (7,-35) rectangle (9,-36);
\fill[stA] (10,-35) rectangle (11,-36);
\fill[stC] (11,-35) rectangle (12,-36);
\fill[stB] (12,-35) rectangle (13,-36);
\fill[stG] (16,-35) rectangle (17,-36);
\fill[stC] (17,-35) rectangle (18,-36);
\fill[stG] (18,-35) rectangle (20,-36);
\fill[stG] (0,-36) rectangle (1,-37);
\fill[stB] (1,-36) rectangle (2,-37);
\fill[stA] (2,-36) rectangle (3,-37);
\fill[stC] (4,-36) rectangle (5,-37);
\fill[stA] (6,-36) rectangle (9,-37);
\fill[stG] (10,-36) rectangle (11,-37);
\fill[stB] (11,-36) rectangle (12,-37);
\fill[stG] (16,-36) rectangle (17,-37);
\fill[stB] (17,-36) rectangle (18,-37);
\fill[stA] (18,-36) rectangle (19,-37);
\fill[stB] (19,-36) rectangle (20,-37);
\fill[stC] (20,-36) rectangle (21,-37);
\fill[stG] (0,-37) rectangle (1,-38);
\fill[stC] (1,-37) rectangle (2,-38);
\fill[stG] (2,-37) rectangle (3,-38);
\fill[stC] (4,-37) rectangle (5,-38);
\fill[stA] (6,-37) rectangle (9,-38);
\fill[stC] (9,-37) rectangle (10,-38);
\fill[stG] (10,-37) rectangle (11,-38);
\fill[stC] (11,-37) rectangle (12,-38);
\fill[stG] (16,-37) rectangle (17,-38);
\fill[stC] (17,-37) rectangle (18,-38);
\fill[stG] (18,-37) rectangle (19,-38);
\fill[stC] (20,-37) rectangle (21,-38);
\fill[stA] (21,-37) rectangle (22,-38);
\fill[stG] (0,-38) rectangle (1,-39);
\fill[stB] (1,-38) rectangle (2,-39);
\fill[stA] (2,-38) rectangle (3,-39);
\fill[stC] (4,-38) rectangle (5,-39);
\fill[stA] (6,-38) rectangle (8,-39);
\fill[stC] (8,-38) rectangle (9,-39);
\fill[stB] (9,-38) rectangle (10,-39);
\fill[stG] (10,-38) rectangle (11,-39);
\fill[stB] (11,-38) rectangle (12,-39);
\fill[stA] (12,-38) rectangle (13,-39);
\fill[stG] (16,-38) rectangle (17,-39);
\fill[stB] (17,-38) rectangle (18,-39);
\fill[stA] (18,-38) rectangle (19,-39);
\fill[stA] (20,-38) rectangle (22,-39);
\fill[stG] (22,-38) rectangle (23,-39);
\fill[stG] (0,-39) rectangle (1,-40);
\fill[stC] (1,-39) rectangle (2,-40);
\fill[stG] (2,-39) rectangle (3,-40);
\fill[stC] (4,-39) rectangle (5,-40);
\fill[stA] (6,-39) rectangle (7,-40);
\fill[stC] (7,-39) rectangle (8,-40);
\fill[stB] (8,-39) rectangle (9,-40);
\fill[stG] (10,-39) rectangle (11,-40);
\fill[stC] (11,-39) rectangle (12,-40);
\fill[stG] (12,-39) rectangle (14,-40);
\fill[stG] (16,-39) rectangle (17,-40);
\fill[stC] (17,-39) rectangle (18,-40);
\fill[stG] (18,-39) rectangle (19,-40);
\fill[stA] (20,-39) rectangle (21,-40);
\fill[stB] (21,-39) rectangle (23,-40);
\fill[stA] (23,-39) rectangle (24,-40);
\fill[stG] (0,-40) rectangle (1,-41);
\fill[stB] (1,-40) rectangle (2,-41);
\fill[stA] (2,-40) rectangle (3,-41);
\fill[stC] (4,-40) rectangle (5,-41);
\fill[stG] (6,-40) rectangle (7,-41);
\fill[stB] (7,-40) rectangle (8,-41);
\fill[stG] (10,-40) rectangle (11,-41);
\fill[stB] (11,-40) rectangle (12,-41);
\fill[stA] (12,-40) rectangle (13,-41);
\fill[stB] (13,-40) rectangle (14,-41);
\fill[stC] (14,-40) rectangle (15,-41);
\fill[stG] (16,-40) rectangle (17,-41);
\fill[stB] (17,-40) rectangle (18,-41);
\fill[stA] (18,-40) rectangle (19,-41);
\fill[stB] (21,-40) rectangle (22,-41);
\fill[stA] (22,-40) rectangle (23,-41);
\fill[stC] (23,-40) rectangle (24,-41);
\fill[stG] (0,-41) rectangle (1,-42);
\fill[stC] (1,-41) rectangle (2,-42);
\fill[stG] (2,-41) rectangle (3,-42);
\fill[stC] (4,-41) rectangle (5,-42);
\fill[stG] (5,-41) rectangle (7,-42);
\fill[stC] (7,-41) rectangle (8,-42);
\fill[stG] (10,-41) rectangle (11,-42);
\fill[stC] (11,-41) rectangle (12,-42);
\fill[stG] (12,-41) rectangle (13,-42);
\fill[stC] (14,-41) rectangle (15,-42);
\fill[stG] (15,-41) rectangle (17,-42);
\fill[stC] (17,-41) rectangle (18,-42);
\fill[stG] (18,-41) rectangle (20,-42);
\fill[stG] (21,-41) rectangle (22,-42);
\fill[stC] (22,-41) rectangle (23,-42);
\fill[stB] (23,-41) rectangle (24,-42);
\fill[stG] (0,-42) rectangle (1,-43);
\fill[stB] (1,-42) rectangle (2,-43);
\fill[stA] (2,-42) rectangle (3,-43);
\fill[stG] (4,-42) rectangle (5,-43);
\fill[stA] (5,-42) rectangle (6,-43);
\fill[stG] (6,-42) rectangle (7,-43);
\fill[stB] (7,-42) rectangle (8,-43);
\fill[stA] (8,-42) rectangle (9,-43);
\fill[stG] (10,-42) rectangle (11,-43);
\fill[stB] (11,-42) rectangle (12,-43);
\fill[stA] (12,-42) rectangle (13,-43);
\fill[stG] (14,-42) rectangle (15,-43);
\fill[stA] (15,-42) rectangle (16,-43);
\fill[stG] (16,-42) rectangle (17,-43);
\fill[stB] (17,-42) rectangle (18,-43);
\fill[stA] (18,-42) rectangle (19,-43);
\fill[stB] (19,-42) rectangle (20,-43);
\fill[stA] (20,-42) rectangle (21,-43);
\fill[stG] (21,-42) rectangle (22,-43);
\fill[stB] (22,-42) rectangle (23,-43);
\fill[stG] (0,-43) rectangle (1,-44);
\fill[stC] (1,-43) rectangle (2,-44);
\fill[stG] (2,-43) rectangle (3,-44);
\fill[stC] (3,-43) rectangle (4,-44);
\fill[stG] (4,-43) rectangle (5,-44);
\fill[stC] (5,-43) rectangle (6,-44);
\fill[stG] (6,-43) rectangle (7,-44);
\fill[stC] (7,-43) rectangle (8,-44);
\fill[stG] (8,-43) rectangle (9,-44);
\fill[stC] (9,-43) rectangle (10,-44);
\fill[stG] (10,-43) rectangle (11,-44);
\fill[stC] (11,-43) rectangle (12,-44);
\fill[stG] (12,-43) rectangle (13,-44);
\fill[stC] (13,-43) rectangle (14,-44);
\fill[stG] (14,-43) rectangle (15,-44);
\fill[stC] (15,-43) rectangle (16,-44);
\fill[stG] (16,-43) rectangle (17,-44);
\fill[stC] (17,-43) rectangle (18,-44);
\fill[stG] (18,-43) rectangle (19,-44);
\fill[stB] (19,-43) rectangle (20,-44);
\fill[stC] (20,-43) rectangle (21,-44);
\fill[stG] (21,-43) rectangle (22,-44);
\fill[stC] (22,-43) rectangle (23,-44);
\fill[stG] (0,-44) rectangle (1,-45);
\fill[stB] (1,-44) rectangle (2,-45);
\fill[stG] (2,-44) rectangle (3,-45);
\fill[stB] (3,-44) rectangle (4,-45);
\fill[stG] (4,-44) rectangle (5,-45);
\fill[stB] (5,-44) rectangle (6,-45);
\fill[stG] (6,-44) rectangle (7,-45);
\fill[stB] (7,-44) rectangle (8,-45);
\fill[stG] (8,-44) rectangle (9,-45);
\fill[stB] (9,-44) rectangle (10,-45);
\fill[stG] (10,-44) rectangle (11,-45);
\fill[stB] (11,-44) rectangle (12,-45);
\fill[stG] (12,-44) rectangle (13,-45);
\fill[stB] (13,-44) rectangle (14,-45);
\fill[stG] (14,-44) rectangle (15,-45);
\fill[stB] (15,-44) rectangle (16,-45);
\fill[stG] (16,-44) rectangle (17,-45);
\fill[stB] (17,-44) rectangle (18,-45);
\fill[stG] (18,-44) rectangle (19,-45);
\fill[stB] (19,-44) rectangle (20,-45);
\fill[stG] (20,-44) rectangle (22,-45);
\fill[stB] (22,-44) rectangle (23,-45);
\fill[stG] (23,-44) rectangle (24,-45);
\fill[stG] (0,-45) rectangle (24,-46);
\fill[stF] (0,-46) rectangle (24,-47);

    \draw[white,line width=2.0pt] \staircase{46}{23} -- (24,-23);
    \draw[ink,densely dashed,line width=0.8pt] \staircase{46}{23} -- (24,-23);
    \draw[grid,line width=0.3pt] (0,0) rectangle (24,-47);
    \node[font=\footnotesize,fill=white,inner sep=1.2pt,rounded corners=1pt] at (17.2,-41.6) {$R_n$};
    \node[font=\small,anchor=south] at (12,0.6) {(a) line, $n=24$};
    \foreach \y/\l in {0/{t=0},23/{t=n-1},46/{t=2n-2}}
      \node[font=\scriptsize,anchor=east] at (-0.5,-\y-0.5) {$\l$};
    \foreach \x/\l in {1/1,24/24}
      \node[font=\scriptsize,anchor=north] at (\x-0.5,-47.3) {$\l$};
  \end{scope}
  % (b) the half-line, cells 1..47
  \begin{scope}[shift={(29,0)}]
    %% generated by figures.py; do not edit
\fill[stG] (0,0) rectangle (1,-1);
\fill[stA] (0,-1) rectangle (1,-2);
\fill[stC] (1,-1) rectangle (2,-2);
\fill[stG] (0,-2) rectangle (1,-3);
\fill[stB] (1,-2) rectangle (2,-3);
\fill[stA] (2,-2) rectangle (3,-3);
\fill[stG] (0,-3) rectangle (1,-4);
\fill[stC] (1,-3) rectangle (2,-4);
\fill[stG] (2,-3) rectangle (4,-4);
\fill[stG] (0,-4) rectangle (1,-5);
\fill[stB] (1,-4) rectangle (2,-5);
\fill[stA] (2,-4) rectangle (3,-5);
\fill[stB] (3,-4) rectangle (4,-5);
\fill[stC] (4,-4) rectangle (5,-5);
\fill[stG] (0,-5) rectangle (1,-6);
\fill[stC] (1,-5) rectangle (2,-6);
\fill[stG] (2,-5) rectangle (3,-6);
\fill[stC] (4,-5) rectangle (5,-6);
\fill[stA] (5,-5) rectangle (6,-6);
\fill[stG] (0,-6) rectangle (1,-7);
\fill[stB] (1,-6) rectangle (2,-7);
\fill[stA] (2,-6) rectangle (3,-7);
\fill[stA] (4,-6) rectangle (6,-7);
\fill[stG] (6,-6) rectangle (7,-7);
\fill[stG] (0,-7) rectangle (1,-8);
\fill[stC] (1,-7) rectangle (2,-8);
\fill[stG] (2,-7) rectangle (3,-8);
\fill[stA] (4,-7) rectangle (5,-8);
\fill[stB] (5,-7) rectangle (7,-8);
\fill[stC] (7,-7) rectangle (8,-8);
\fill[stG] (0,-8) rectangle (1,-9);
\fill[stB] (1,-8) rectangle (2,-9);
\fill[stA] (2,-8) rectangle (3,-9);
\fill[stB] (5,-8) rectangle (6,-9);
\fill[stC] (6,-8) rectangle (8,-9);
\fill[stA] (8,-8) rectangle (9,-9);
\fill[stG] (0,-9) rectangle (1,-10);
\fill[stC] (1,-9) rectangle (2,-10);
\fill[stG] (2,-9) rectangle (4,-10);
\fill[stC] (6,-9) rectangle (7,-10);
\fill[stA] (7,-9) rectangle (9,-10);
\fill[stG] (9,-9) rectangle (10,-10);
\fill[stG] (0,-10) rectangle (1,-11);
\fill[stB] (1,-10) rectangle (2,-11);
\fill[stA] (2,-10) rectangle (3,-11);
\fill[stB] (3,-10) rectangle (4,-11);
\fill[stC] (4,-10) rectangle (5,-11);
\fill[stA] (6,-10) rectangle (8,-11);
\fill[stB] (8,-10) rectangle (10,-11);
\fill[stC] (10,-10) rectangle (11,-11);
\fill[stG] (0,-11) rectangle (1,-12);
\fill[stC] (1,-11) rectangle (2,-12);
\fill[stG] (2,-11) rectangle (3,-12);
\fill[stC] (4,-11) rectangle (5,-12);
\fill[stA] (6,-11) rectangle (7,-12);
\fill[stB] (7,-11) rectangle (9,-12);
\fill[stC] (9,-11) rectangle (11,-12);
\fill[stA] (11,-11) rectangle (12,-12);
\fill[stG] (0,-12) rectangle (1,-13);
\fill[stB] (1,-12) rectangle (2,-13);
\fill[stA] (2,-12) rectangle (3,-13);
\fill[stC] (4,-12) rectangle (5,-13);
\fill[stB] (7,-12) rectangle (8,-13);
\fill[stC] (8,-12) rectangle (10,-13);
\fill[stA] (10,-12) rectangle (12,-13);
\fill[stG] (12,-12) rectangle (13,-13);
\fill[stG] (0,-13) rectangle (1,-14);
\fill[stC] (1,-13) rectangle (2,-14);
\fill[stG] (2,-13) rectangle (3,-14);
\fill[stC] (4,-13) rectangle (5,-14);
\fill[stA] (5,-13) rectangle (6,-14);
\fill[stC] (8,-13) rectangle (9,-14);
\fill[stA] (9,-13) rectangle (11,-14);
\fill[stB] (11,-13) rectangle (13,-14);
\fill[stC] (13,-13) rectangle (14,-14);
\fill[stG] (0,-14) rectangle (1,-15);
\fill[stB] (1,-14) rectangle (2,-15);
\fill[stA] (2,-14) rectangle (3,-15);
\fill[stA] (4,-14) rectangle (6,-15);
\fill[stG] (6,-14) rectangle (7,-15);
\fill[stA] (8,-14) rectangle (10,-15);
\fill[stB] (10,-14) rectangle (12,-15);
\fill[stC] (12,-14) rectangle (14,-15);
\fill[stA] (14,-14) rectangle (15,-15);
\fill[stG] (0,-15) rectangle (1,-16);
\fill[stC] (1,-15) rectangle (2,-16);
\fill[stG] (2,-15) rectangle (3,-16);
\fill[stA] (4,-15) rectangle (5,-16);
\fill[stB] (5,-15) rectangle (7,-16);
\fill[stA] (8,-15) rectangle (9,-16);
\fill[stB] (9,-15) rectangle (11,-16);
\fill[stC] (11,-15) rectangle (13,-16);
\fill[stA] (13,-15) rectangle (15,-16);
\fill[stG] (15,-15) rectangle (16,-16);
\fill[stG] (0,-16) rectangle (1,-17);
\fill[stB] (1,-16) rectangle (2,-17);
\fill[stA] (2,-16) rectangle (3,-17);
\fill[stB] (5,-16) rectangle (6,-17);
\fill[stG] (6,-16) rectangle (7,-17);
\fill[stB] (9,-16) rectangle (10,-17);
\fill[stC] (10,-16) rectangle (12,-17);
\fill[stA] (12,-16) rectangle (14,-17);
\fill[stB] (14,-16) rectangle (16,-17);
\fill[stC] (16,-16) rectangle (17,-17);
\fill[stG] (0,-17) rectangle (1,-18);
\fill[stC] (1,-17) rectangle (2,-18);
\fill[stG] (2,-17) rectangle (4,-18);
\fill[stB] (5,-17) rectangle (7,-18);
\fill[stC] (7,-17) rectangle (8,-18);
\fill[stC] (10,-17) rectangle (11,-18);
\fill[stA] (11,-17) rectangle (13,-18);
\fill[stB] (13,-17) rectangle (15,-18);
\fill[stC] (15,-17) rectangle (17,-18);
\fill[stA] (17,-17) rectangle (18,-18);
\fill[stG] (0,-18) rectangle (1,-19);
\fill[stB] (1,-18) rectangle (2,-19);
\fill[stA] (2,-18) rectangle (3,-19);
\fill[stB] (3,-18) rectangle (4,-19);
\fill[stB] (5,-18) rectangle (6,-19);
\fill[stC] (6,-18) rectangle (8,-19);
\fill[stA] (8,-18) rectangle (9,-19);
\fill[stA] (10,-18) rectangle (12,-19);
\fill[stB] (12,-18) rectangle (14,-19);
\fill[stC] (14,-18) rectangle (16,-19);
\fill[stA] (16,-18) rectangle (18,-19);
\fill[stG] (18,-18) rectangle (19,-19);
\fill[stG] (0,-19) rectangle (1,-20);
\fill[stC] (1,-19) rectangle (2,-20);
\fill[stG] (2,-19) rectangle (4,-20);
\fill[stC] (6,-19) rectangle (7,-20);
\fill[stA] (7,-19) rectangle (9,-20);
\fill[stA] (10,-19) rectangle (11,-20);
\fill[stB] (11,-19) rectangle (13,-20);
\fill[stC] (13,-19) rectangle (15,-20);
\fill[stA] (15,-19) rectangle (17,-20);
\fill[stB] (17,-19) rectangle (19,-20);
\fill[stC] (19,-19) rectangle (20,-20);
\fill[stG] (0,-20) rectangle (1,-21);
\fill[stB] (1,-20) rectangle (2,-21);
\fill[stA] (2,-20) rectangle (3,-21);
\fill[stB] (3,-20) rectangle (4,-21);
\fill[stC] (4,-20) rectangle (5,-21);
\fill[stA] (6,-20) rectangle (9,-21);
\fill[stB] (11,-20) rectangle (12,-21);
\fill[stC] (12,-20) rectangle (14,-21);
\fill[stA] (14,-20) rectangle (16,-21);
\fill[stB] (16,-20) rectangle (18,-21);
\fill[stC] (18,-20) rectangle (20,-21);
\fill[stA] (20,-20) rectangle (21,-21);
\fill[stG] (0,-21) rectangle (1,-22);
\fill[stC] (1,-21) rectangle (2,-22);
\fill[stG] (2,-21) rectangle (3,-22);
\fill[stC] (4,-21) rectangle (5,-22);
\fill[stA] (6,-21) rectangle (9,-22);
\fill[stG] (9,-21) rectangle (10,-22);
\fill[stC] (12,-21) rectangle (13,-22);
\fill[stA] (13,-21) rectangle (15,-22);
\fill[stB] (15,-21) rectangle (17,-22);
\fill[stC] (17,-21) rectangle (19,-22);
\fill[stA] (19,-21) rectangle (21,-22);
\fill[stG] (21,-21) rectangle (22,-22);
\fill[stG] (0,-22) rectangle (1,-23);
\fill[stB] (1,-22) rectangle (2,-23);
\fill[stA] (2,-22) rectangle (3,-23);
\fill[stC] (4,-22) rectangle (5,-23);
\fill[stA] (6,-22) rectangle (8,-23);
\fill[stB] (8,-22) rectangle (10,-23);
\fill[stC] (10,-22) rectangle (11,-23);
\fill[stA] (12,-22) rectangle (14,-23);
\fill[stB] (14,-22) rectangle (16,-23);
\fill[stC] (16,-22) rectangle (18,-23);
\fill[stA] (18,-22) rectangle (20,-23);
\fill[stB] (20,-22) rectangle (22,-23);
\fill[stC] (22,-22) rectangle (23,-23);
\fill[stG] (0,-23) rectangle (1,-24);
\fill[stC] (1,-23) rectangle (2,-24);
\fill[stG] (2,-23) rectangle (3,-24);
\fill[stC] (4,-23) rectangle (5,-24);
\fill[stA] (6,-23) rectangle (7,-24);
\fill[stB] (7,-23) rectangle (9,-24);
\fill[stC] (9,-23) rectangle (11,-24);
\fill[stA] (12,-23) rectangle (13,-24);
\fill[stB] (13,-23) rectangle (15,-24);
\fill[stC] (15,-23) rectangle (17,-24);
\fill[stA] (17,-23) rectangle (19,-24);
\fill[stB] (19,-23) rectangle (21,-24);
\fill[stC] (21,-23) rectangle (23,-24);
\fill[stA] (23,-23) rectangle (24,-24);
\fill[stG] (0,-24) rectangle (1,-25);
\fill[stB] (1,-24) rectangle (2,-25);
\fill[stA] (2,-24) rectangle (3,-25);
\fill[stC] (4,-24) rectangle (5,-25);
\fill[stB] (7,-24) rectangle (8,-25);
\fill[stC] (8,-24) rectangle (11,-25);
\fill[stB] (13,-24) rectangle (14,-25);
\fill[stC] (14,-24) rectangle (16,-25);
\fill[stA] (16,-24) rectangle (18,-25);
\fill[stB] (18,-24) rectangle (20,-25);
\fill[stC] (20,-24) rectangle (22,-25);
\fill[stA] (22,-24) rectangle (24,-25);
\fill[stG] (24,-24) rectangle (25,-25);
\fill[stG] (0,-25) rectangle (1,-26);
\fill[stC] (1,-25) rectangle (2,-26);
\fill[stG] (2,-25) rectangle (3,-26);
\fill[stC] (4,-25) rectangle (5,-26);
\fill[stA] (5,-25) rectangle (6,-26);
\fill[stC] (8,-25) rectangle (11,-26);
\fill[stA] (11,-25) rectangle (12,-26);
\fill[stC] (14,-25) rectangle (15,-26);
\fill[stA] (15,-25) rectangle (17,-26);
\fill[stB] (17,-25) rectangle (19,-26);
\fill[stC] (19,-25) rectangle (21,-26);
\fill[stA] (21,-25) rectangle (23,-26);
\fill[stB] (23,-25) rectangle (25,-26);
\fill[stC] (25,-25) rectangle (26,-26);
\fill[stG] (0,-26) rectangle (1,-27);
\fill[stB] (1,-26) rectangle (2,-27);
\fill[stA] (2,-26) rectangle (3,-27);
\fill[stA] (4,-26) rectangle (6,-27);
\fill[stG] (6,-26) rectangle (7,-27);
\fill[stC] (8,-26) rectangle (10,-27);
\fill[stA] (10,-26) rectangle (12,-27);
\fill[stG] (12,-26) rectangle (13,-27);
\fill[stA] (14,-26) rectangle (16,-27);
\fill[stB] (16,-26) rectangle (18,-27);
\fill[stC] (18,-26) rectangle (20,-27);
\fill[stA] (20,-26) rectangle (22,-27);
\fill[stB] (22,-26) rectangle (24,-27);
\fill[stC] (24,-26) rectangle (26,-27);
\fill[stA] (26,-26) rectangle (27,-27);
\fill[stG] (0,-27) rectangle (1,-28);
\fill[stC] (1,-27) rectangle (2,-28);
\fill[stG] (2,-27) rectangle (3,-28);
\fill[stA] (4,-27) rectangle (5,-28);
\fill[stB] (5,-27) rectangle (7,-28);
\fill[stC] (8,-27) rectangle (9,-28);
\fill[stA] (9,-27) rectangle (11,-28);
\fill[stB] (11,-27) rectangle (13,-28);
\fill[stA] (14,-27) rectangle (15,-28);
\fill[stB] (15,-27) rectangle (17,-28);
\fill[stC] (17,-27) rectangle (19,-28);
\fill[stA] (19,-27) rectangle (21,-28);
\fill[stB] (21,-27) rectangle (23,-28);
\fill[stC] (23,-27) rectangle (25,-28);
\fill[stA] (25,-27) rectangle (27,-28);
\fill[stG] (27,-27) rectangle (28,-28);
\fill[stG] (0,-28) rectangle (1,-29);
\fill[stB] (1,-28) rectangle (2,-29);
\fill[stA] (2,-28) rectangle (3,-29);
\fill[stB] (5,-28) rectangle (6,-29);
\fill[stG] (6,-28) rectangle (7,-29);
\fill[stA] (8,-28) rectangle (10,-29);
\fill[stB] (10,-28) rectangle (12,-29);
\fill[stG] (12,-28) rectangle (13,-29);
\fill[stB] (15,-28) rectangle (16,-29);
\fill[stC] (16,-28) rectangle (18,-29);
\fill[stA] (18,-28) rectangle (20,-29);
\fill[stB] (20,-28) rectangle (22,-29);
\fill[stC] (22,-28) rectangle (24,-29);
\fill[stA] (24,-28) rectangle (26,-29);
\fill[stB] (26,-28) rectangle (28,-29);
\fill[stC] (28,-28) rectangle (29,-29);
\fill[stG] (0,-29) rectangle (1,-30);
\fill[stC] (1,-29) rectangle (2,-30);
\fill[stG] (2,-29) rectangle (4,-30);
\fill[stB] (5,-29) rectangle (7,-30);
\fill[stA] (8,-29) rectangle (9,-30);
\fill[stB] (9,-29) rectangle (13,-30);
\fill[stC] (13,-29) rectangle (14,-30);
\fill[stC] (16,-29) rectangle (17,-30);
\fill[stA] (17,-29) rectangle (19,-30);
\fill[stB] (19,-29) rectangle (21,-30);
\fill[stC] (21,-29) rectangle (23,-30);
\fill[stA] (23,-29) rectangle (25,-30);
\fill[stB] (25,-29) rectangle (27,-30);
\fill[stC] (27,-29) rectangle (29,-30);
\fill[stA] (29,-29) rectangle (30,-30);
\fill[stG] (0,-30) rectangle (1,-31);
\fill[stB] (1,-30) rectangle (2,-31);
\fill[stA] (2,-30) rectangle (3,-31);
\fill[stB] (3,-30) rectangle (4,-31);
\fill[stB] (5,-30) rectangle (6,-31);
\fill[stG] (6,-30) rectangle (7,-31);
\fill[stB] (9,-30) rectangle (12,-31);
\fill[stC] (12,-30) rectangle (14,-31);
\fill[stA] (14,-30) rectangle (15,-31);
\fill[stA] (16,-30) rectangle (18,-31);
\fill[stB] (18,-30) rectangle (20,-31);
\fill[stC] (20,-30) rectangle (22,-31);
\fill[stA] (22,-30) rectangle (24,-31);
\fill[stB] (24,-30) rectangle (26,-31);
\fill[stC] (26,-30) rectangle (28,-31);
\fill[stA] (28,-30) rectangle (30,-31);
\fill[stG] (30,-30) rectangle (31,-31);
\fill[stG] (0,-31) rectangle (1,-32);
\fill[stC] (1,-31) rectangle (2,-32);
\fill[stG] (2,-31) rectangle (4,-32);
\fill[stB] (5,-31) rectangle (7,-32);
\fill[stC] (7,-31) rectangle (8,-32);
\fill[stB] (9,-31) rectangle (11,-32);
\fill[stC] (11,-31) rectangle (13,-32);
\fill[stA] (13,-31) rectangle (15,-32);
\fill[stA] (16,-31) rectangle (17,-32);
\fill[stB] (17,-31) rectangle (19,-32);
\fill[stC] (19,-31) rectangle (21,-32);
\fill[stA] (21,-31) rectangle (23,-32);
\fill[stB] (23,-31) rectangle (25,-32);
\fill[stC] (25,-31) rectangle (27,-32);
\fill[stA] (27,-31) rectangle (29,-32);
\fill[stB] (29,-31) rectangle (31,-32);
\fill[stC] (31,-31) rectangle (32,-32);
\fill[stG] (0,-32) rectangle (1,-33);
\fill[stB] (1,-32) rectangle (2,-33);
\fill[stA] (2,-32) rectangle (3,-33);
\fill[stB] (3,-32) rectangle (4,-33);
\fill[stB] (5,-32) rectangle (6,-33);
\fill[stC] (6,-32) rectangle (8,-33);
\fill[stB] (9,-32) rectangle (10,-33);
\fill[stC] (10,-32) rectangle (12,-33);
\fill[stA] (12,-32) rectangle (15,-33);
\fill[stB] (17,-32) rectangle (18,-33);
\fill[stC] (18,-32) rectangle (20,-33);
\fill[stA] (20,-32) rectangle (22,-33);
\fill[stB] (22,-32) rectangle (24,-33);
\fill[stC] (24,-32) rectangle (26,-33);
\fill[stA] (26,-32) rectangle (28,-33);
\fill[stB] (28,-32) rectangle (30,-33);
\fill[stC] (30,-32) rectangle (32,-33);
\fill[stA] (32,-32) rectangle (33,-33);
\fill[stG] (0,-33) rectangle (1,-34);
\fill[stC] (1,-33) rectangle (2,-34);
\fill[stG] (2,-33) rectangle (4,-34);
\fill[stC] (6,-33) rectangle (8,-34);
\fill[stC] (10,-33) rectangle (11,-34);
\fill[stA] (11,-33) rectangle (15,-34);
\fill[stG] (15,-33) rectangle (16,-34);
\fill[stC] (18,-33) rectangle (19,-34);
\fill[stA] (19,-33) rectangle (21,-34);
\fill[stB] (21,-33) rectangle (23,-34);
\fill[stC] (23,-33) rectangle (25,-34);
\fill[stA] (25,-33) rectangle (27,-34);
\fill[stB] (27,-33) rectangle (29,-34);
\fill[stC] (29,-33) rectangle (31,-34);
\fill[stA] (31,-33) rectangle (33,-34);
\fill[stG] (33,-33) rectangle (34,-34);
\fill[stG] (0,-34) rectangle (1,-35);
\fill[stB] (1,-34) rectangle (2,-35);
\fill[stA] (2,-34) rectangle (3,-35);
\fill[stB] (3,-34) rectangle (4,-35);
\fill[stC] (4,-34) rectangle (5,-35);
\fill[stC] (6,-34) rectangle (8,-35);
\fill[stA] (8,-34) rectangle (9,-35);
\fill[stA] (10,-34) rectangle (14,-35);
\fill[stB] (14,-34) rectangle (16,-35);
\fill[stC] (16,-34) rectangle (17,-35);
\fill[stA] (18,-34) rectangle (20,-35);
\fill[stB] (20,-34) rectangle (22,-35);
\fill[stC] (22,-34) rectangle (24,-35);
\fill[stA] (24,-34) rectangle (26,-35);
\fill[stB] (26,-34) rectangle (28,-35);
\fill[stC] (28,-34) rectangle (30,-35);
\fill[stA] (30,-34) rectangle (32,-35);
\fill[stB] (32,-34) rectangle (34,-35);
\fill[stC] (34,-34) rectangle (35,-35);
\fill[stG] (0,-35) rectangle (1,-36);
\fill[stC] (1,-35) rectangle (2,-36);
\fill[stG] (2,-35) rectangle (3,-36);
\fill[stC] (4,-35) rectangle (5,-36);
\fill[stC] (6,-35) rectangle (7,-36);
\fill[stA] (7,-35) rectangle (9,-36);
\fill[stA] (10,-35) rectangle (13,-36);
\fill[stB] (13,-35) rectangle (15,-36);
\fill[stC] (15,-35) rectangle (17,-36);
\fill[stA] (18,-35) rectangle (19,-36);
\fill[stB] (19,-35) rectangle (21,-36);
\fill[stC] (21,-35) rectangle (23,-36);
\fill[stA] (23,-35) rectangle (25,-36);
\fill[stB] (25,-35) rectangle (27,-36);
\fill[stC] (27,-35) rectangle (29,-36);
\fill[stA] (29,-35) rectangle (31,-36);
\fill[stB] (31,-35) rectangle (33,-36);
\fill[stC] (33,-35) rectangle (35,-36);
\fill[stA] (35,-35) rectangle (36,-36);
\fill[stG] (0,-36) rectangle (1,-37);
\fill[stB] (1,-36) rectangle (2,-37);
\fill[stA] (2,-36) rectangle (3,-37);
\fill[stC] (4,-36) rectangle (5,-37);
\fill[stA] (6,-36) rectangle (9,-37);
\fill[stA] (10,-36) rectangle (12,-37);
\fill[stB] (12,-36) rectangle (14,-37);
\fill[stC] (14,-36) rectangle (17,-37);
\fill[stB] (19,-36) rectangle (20,-37);
\fill[stC] (20,-36) rectangle (22,-37);
\fill[stA] (22,-36) rectangle (24,-37);
\fill[stB] (24,-36) rectangle (26,-37);
\fill[stC] (26,-36) rectangle (28,-37);
\fill[stA] (28,-36) rectangle (30,-37);
\fill[stB] (30,-36) rectangle (32,-37);
\fill[stC] (32,-36) rectangle (34,-37);
\fill[stA] (34,-36) rectangle (36,-37);
\fill[stG] (36,-36) rectangle (37,-37);
\fill[stG] (0,-37) rectangle (1,-38);
\fill[stC] (1,-37) rectangle (2,-38);
\fill[stG] (2,-37) rectangle (3,-38);
\fill[stC] (4,-37) rectangle (5,-38);
\fill[stA] (6,-37) rectangle (9,-38);
\fill[stA] (10,-37) rectangle (11,-38);
\fill[stB] (11,-37) rectangle (13,-38);
\fill[stC] (13,-37) rectangle (17,-38);
\fill[stA] (17,-37) rectangle (18,-38);
\fill[stC] (20,-37) rectangle (21,-38);
\fill[stA] (21,-37) rectangle (23,-38);
\fill[stB] (23,-37) rectangle (25,-38);
\fill[stC] (25,-37) rectangle (27,-38);
\fill[stA] (27,-37) rectangle (29,-38);
\fill[stB] (29,-37) rectangle (31,-38);
\fill[stC] (31,-37) rectangle (33,-38);
\fill[stA] (33,-37) rectangle (35,-38);
\fill[stB] (35,-37) rectangle (37,-38);
\fill[stC] (37,-37) rectangle (38,-38);
\fill[stG] (0,-38) rectangle (1,-39);
\fill[stB] (1,-38) rectangle (2,-39);
\fill[stA] (2,-38) rectangle (3,-39);
\fill[stC] (4,-38) rectangle (5,-39);
\fill[stA] (6,-38) rectangle (9,-39);
\fill[stB] (11,-38) rectangle (12,-39);
\fill[stC] (12,-38) rectangle (16,-39);
\fill[stA] (16,-38) rectangle (18,-39);
\fill[stG] (18,-38) rectangle (19,-39);
\fill[stA] (20,-38) rectangle (22,-39);
\fill[stB] (22,-38) rectangle (24,-39);
\fill[stC] (24,-38) rectangle (26,-39);
\fill[stA] (26,-38) rectangle (28,-39);
\fill[stB] (28,-38) rectangle (30,-39);
\fill[stC] (30,-38) rectangle (32,-39);
\fill[stA] (32,-38) rectangle (34,-39);
\fill[stB] (34,-38) rectangle (36,-39);
\fill[stC] (36,-38) rectangle (38,-39);
\fill[stA] (38,-38) rectangle (39,-39);
\fill[stG] (0,-39) rectangle (1,-40);
\fill[stC] (1,-39) rectangle (2,-40);
\fill[stG] (2,-39) rectangle (3,-40);
\fill[stC] (4,-39) rectangle (5,-40);
\fill[stA] (6,-39) rectangle (9,-40);
\fill[stG] (9,-39) rectangle (10,-40);
\fill[stC] (12,-39) rectangle (15,-40);
\fill[stA] (15,-39) rectangle (17,-40);
\fill[stB] (17,-39) rectangle (19,-40);
\fill[stA] (20,-39) rectangle (21,-40);
\fill[stB] (21,-39) rectangle (23,-40);
\fill[stC] (23,-39) rectangle (25,-40);
\fill[stA] (25,-39) rectangle (27,-40);
\fill[stB] (27,-39) rectangle (29,-40);
\fill[stC] (29,-39) rectangle (31,-40);
\fill[stA] (31,-39) rectangle (33,-40);
\fill[stB] (33,-39) rectangle (35,-40);
\fill[stC] (35,-39) rectangle (37,-40);
\fill[stA] (37,-39) rectangle (39,-40);
\fill[stG] (39,-39) rectangle (40,-40);
\fill[stG] (0,-40) rectangle (1,-41);
\fill[stB] (1,-40) rectangle (2,-41);
\fill[stA] (2,-40) rectangle (3,-41);
\fill[stC] (4,-40) rectangle (5,-41);
\fill[stA] (6,-40) rectangle (8,-41);
\fill[stB] (8,-40) rectangle (10,-41);
\fill[stC] (10,-40) rectangle (11,-41);
\fill[stC] (12,-40) rectangle (14,-41);
\fill[stA] (14,-40) rectangle (16,-41);
\fill[stB] (16,-40) rectangle (18,-41);
\fill[stG] (18,-40) rectangle (19,-41);
\fill[stB] (21,-40) rectangle (22,-41);
\fill[stC] (22,-40) rectangle (24,-41);
\fill[stA] (24,-40) rectangle (26,-41);
\fill[stB] (26,-40) rectangle (28,-41);
\fill[stC] (28,-40) rectangle (30,-41);
\fill[stA] (30,-40) rectangle (32,-41);
\fill[stB] (32,-40) rectangle (34,-41);
\fill[stC] (34,-40) rectangle (36,-41);
\fill[stA] (36,-40) rectangle (38,-41);
\fill[stB] (38,-40) rectangle (40,-41);
\fill[stC] (40,-40) rectangle (41,-41);
\fill[stG] (0,-41) rectangle (1,-42);
\fill[stC] (1,-41) rectangle (2,-42);
\fill[stG] (2,-41) rectangle (3,-42);
\fill[stC] (4,-41) rectangle (5,-42);
\fill[stA] (6,-41) rectangle (7,-42);
\fill[stB] (7,-41) rectangle (9,-42);
\fill[stC] (9,-41) rectangle (11,-42);
\fill[stC] (12,-41) rectangle (13,-42);
\fill[stA] (13,-41) rectangle (15,-42);
\fill[stB] (15,-41) rectangle (19,-42);
\fill[stC] (19,-41) rectangle (20,-42);
\fill[stC] (22,-41) rectangle (23,-42);
\fill[stA] (23,-41) rectangle (25,-42);
\fill[stB] (25,-41) rectangle (27,-42);
\fill[stC] (27,-41) rectangle (29,-42);
\fill[stA] (29,-41) rectangle (31,-42);
\fill[stB] (31,-41) rectangle (33,-42);
\fill[stC] (33,-41) rectangle (35,-42);
\fill[stA] (35,-41) rectangle (37,-42);
\fill[stB] (37,-41) rectangle (39,-42);
\fill[stC] (39,-41) rectangle (41,-42);
\fill[stA] (41,-41) rectangle (42,-42);
\fill[stG] (0,-42) rectangle (1,-43);
\fill[stB] (1,-42) rectangle (2,-43);
\fill[stA] (2,-42) rectangle (3,-43);
\fill[stC] (4,-42) rectangle (5,-43);
\fill[stB] (7,-42) rectangle (8,-43);
\fill[stC] (8,-42) rectangle (11,-43);
\fill[stA] (12,-42) rectangle (14,-43);
\fill[stB] (14,-42) rectangle (18,-43);
\fill[stC] (18,-42) rectangle (20,-43);
\fill[stA] (20,-42) rectangle (21,-43);
\fill[stA] (22,-42) rectangle (24,-43);
\fill[stB] (24,-42) rectangle (26,-43);
\fill[stC] (26,-42) rectangle (28,-43);
\fill[stA] (28,-42) rectangle (30,-43);
\fill[stB] (30,-42) rectangle (32,-43);
\fill[stC] (32,-42) rectangle (34,-43);
\fill[stA] (34,-42) rectangle (36,-43);
\fill[stB] (36,-42) rectangle (38,-43);
\fill[stC] (38,-42) rectangle (40,-43);
\fill[stA] (40,-42) rectangle (42,-43);
\fill[stG] (42,-42) rectangle (43,-43);
\fill[stG] (0,-43) rectangle (1,-44);
\fill[stC] (1,-43) rectangle (2,-44);
\fill[stG] (2,-43) rectangle (3,-44);
\fill[stC] (4,-43) rectangle (5,-44);
\fill[stA] (5,-43) rectangle (6,-44);
\fill[stC] (8,-43) rectangle (11,-44);
\fill[stA] (12,-43) rectangle (13,-44);
\fill[stB] (13,-43) rectangle (17,-44);
\fill[stC] (17,-43) rectangle (19,-44);
\fill[stA] (19,-43) rectangle (21,-44);
\fill[stA] (22,-43) rectangle (23,-44);
\fill[stB] (23,-43) rectangle (25,-44);
\fill[stC] (25,-43) rectangle (27,-44);
\fill[stA] (27,-43) rectangle (29,-44);
\fill[stB] (29,-43) rectangle (31,-44);
\fill[stC] (31,-43) rectangle (33,-44);
\fill[stA] (33,-43) rectangle (35,-44);
\fill[stB] (35,-43) rectangle (37,-44);
\fill[stC] (37,-43) rectangle (39,-44);
\fill[stA] (39,-43) rectangle (41,-44);
\fill[stB] (41,-43) rectangle (43,-44);
\fill[stC] (43,-43) rectangle (44,-44);
\fill[stG] (0,-44) rectangle (1,-45);
\fill[stB] (1,-44) rectangle (2,-45);
\fill[stA] (2,-44) rectangle (3,-45);
\fill[stA] (4,-44) rectangle (6,-45);
\fill[stG] (6,-44) rectangle (7,-45);
\fill[stC] (8,-44) rectangle (11,-45);
\fill[stB] (13,-44) rectangle (16,-45);
\fill[stC] (16,-44) rectangle (18,-45);
\fill[stA] (18,-44) rectangle (21,-45);
\fill[stB] (23,-44) rectangle (24,-45);
\fill[stC] (24,-44) rectangle (26,-45);
\fill[stA] (26,-44) rectangle (28,-45);
\fill[stB] (28,-44) rectangle (30,-45);
\fill[stC] (30,-44) rectangle (32,-45);
\fill[stA] (32,-44) rectangle (34,-45);
\fill[stB] (34,-44) rectangle (36,-45);
\fill[stC] (36,-44) rectangle (38,-45);
\fill[stA] (38,-44) rectangle (40,-45);
\fill[stB] (40,-44) rectangle (42,-45);
\fill[stC] (42,-44) rectangle (44,-45);
\fill[stA] (44,-44) rectangle (45,-45);
\fill[stG] (0,-45) rectangle (1,-46);
\fill[stC] (1,-45) rectangle (2,-46);
\fill[stG] (2,-45) rectangle (3,-46);
\fill[stA] (4,-45) rectangle (5,-46);
\fill[stB] (5,-45) rectangle (7,-46);
\fill[stC] (8,-45) rectangle (11,-46);
\fill[stA] (11,-45) rectangle (12,-46);
\fill[stB] (13,-45) rectangle (15,-46);
\fill[stC] (15,-45) rectangle (17,-46);
\fill[stA] (17,-45) rectangle (21,-46);
\fill[stG] (21,-45) rectangle (22,-46);
\fill[stC] (24,-45) rectangle (25,-46);
\fill[stA] (25,-45) rectangle (27,-46);
\fill[stB] (27,-45) rectangle (29,-46);
\fill[stC] (29,-45) rectangle (31,-46);
\fill[stA] (31,-45) rectangle (33,-46);
\fill[stB] (33,-45) rectangle (35,-46);
\fill[stC] (35,-45) rectangle (37,-46);
\fill[stA] (37,-45) rectangle (39,-46);
\fill[stB] (39,-45) rectangle (41,-46);
\fill[stC] (41,-45) rectangle (43,-46);
\fill[stA] (43,-45) rectangle (45,-46);
\fill[stG] (45,-45) rectangle (46,-46);
\fill[stG] (0,-46) rectangle (1,-47);
\fill[stB] (1,-46) rectangle (2,-47);
\fill[stA] (2,-46) rectangle (3,-47);
\fill[stB] (5,-46) rectangle (6,-47);
\fill[stG] (6,-46) rectangle (7,-47);
\fill[stC] (8,-46) rectangle (10,-47);
\fill[stA] (10,-46) rectangle (12,-47);
\fill[stB] (13,-46) rectangle (14,-47);
\fill[stC] (14,-46) rectangle (16,-47);
\fill[stA] (16,-46) rectangle (20,-47);
\fill[stB] (20,-46) rectangle (22,-47);
\fill[stC] (22,-46) rectangle (23,-47);
\fill[stA] (24,-46) rectangle (26,-47);
\fill[stB] (26,-46) rectangle (28,-47);
\fill[stC] (28,-46) rectangle (30,-47);
\fill[stA] (30,-46) rectangle (32,-47);
\fill[stB] (32,-46) rectangle (34,-47);
\fill[stC] (34,-46) rectangle (36,-47);
\fill[stA] (36,-46) rectangle (38,-47);
\fill[stB] (38,-46) rectangle (40,-47);
\fill[stC] (40,-46) rectangle (42,-47);
\fill[stA] (42,-46) rectangle (44,-47);
\fill[stB] (44,-46) rectangle (46,-47);
\fill[stC] (46,-46) rectangle (47,-47);

    \draw[ink,densely dotted,line width=0.9pt] (24,0) -- (24,-47);
    \draw[white,line width=2.0pt] \staircase{46}{46};
    \draw[ink,densely dashed,line width=0.8pt] \staircase{46}{46};
    \draw[grid,line width=0.3pt] (0,0) rectangle (47,-47);
    \node[font=\scriptsize,fill=white,inner sep=1pt,rotate=45] at (36.3,-13.9) {$t+i=2n-2$};
    \node[font=\small,anchor=south] at (23.5,0.6) {(b) half-line};
    \foreach \x/\l in {1/1,24/24,47/47}
      \node[font=\scriptsize,anchor=north] at (\x-0.5,-47.3) {$\l$};
  \end{scope}
  % (c) where the line of length 24 differs from the half-line
  \begin{scope}[shift={(81,0)}]
    %% generated by figures.py; do not edit
\fill[diffSame] (0,0) rectangle (1,-1);
\fill[diffSame] (0,-1) rectangle (2,-2);
\fill[diffSame] (0,-2) rectangle (3,-3);
\fill[diffSame] (0,-3) rectangle (4,-4);
\fill[diffSame] (0,-4) rectangle (5,-5);
\fill[diffSame] (0,-5) rectangle (6,-6);
\fill[diffSame] (0,-6) rectangle (7,-7);
\fill[diffSame] (0,-7) rectangle (8,-8);
\fill[diffSame] (0,-8) rectangle (9,-9);
\fill[diffSame] (0,-9) rectangle (10,-10);
\fill[diffSame] (0,-10) rectangle (11,-11);
\fill[diffSame] (0,-11) rectangle (12,-12);
\fill[diffSame] (0,-12) rectangle (13,-13);
\fill[diffSame] (0,-13) rectangle (14,-14);
\fill[diffSame] (0,-14) rectangle (15,-15);
\fill[diffSame] (0,-15) rectangle (16,-16);
\fill[diffSame] (0,-16) rectangle (17,-17);
\fill[diffSame] (0,-17) rectangle (18,-18);
\fill[diffSame] (0,-18) rectangle (19,-19);
\fill[diffSame] (0,-19) rectangle (20,-20);
\fill[diffSame] (0,-20) rectangle (21,-21);
\fill[diffSame] (0,-21) rectangle (22,-22);
\fill[diffSame] (0,-22) rectangle (23,-23);
\fill[diffSame] (0,-23) rectangle (23,-24);
\fill[diffDiff] (23,-23) rectangle (24,-24);
\fill[diffSame] (0,-24) rectangle (22,-25);
\fill[diffDiff] (22,-24) rectangle (23,-25);
\fill[diffSame] (23,-24) rectangle (24,-25);
\fill[diffSame] (0,-25) rectangle (21,-26);
\fill[diffDiff] (21,-25) rectangle (22,-26);
\fill[diffSame] (22,-25) rectangle (23,-26);
\fill[diffDiff] (23,-25) rectangle (24,-26);
\fill[diffSame] (0,-26) rectangle (20,-27);
\fill[diffDiff] (20,-26) rectangle (21,-27);
\fill[diffSame] (21,-26) rectangle (22,-27);
\fill[diffDiff] (22,-26) rectangle (23,-27);
\fill[diffSame] (23,-26) rectangle (24,-27);
\fill[diffSame] (0,-27) rectangle (19,-28);
\fill[diffDiff] (19,-27) rectangle (20,-28);
\fill[diffSame] (20,-27) rectangle (21,-28);
\fill[diffDiff] (21,-27) rectangle (22,-28);
\fill[diffSame] (22,-27) rectangle (23,-28);
\fill[diffDiff] (23,-27) rectangle (24,-28);
\fill[diffSame] (0,-28) rectangle (18,-29);
\fill[diffDiff] (18,-28) rectangle (19,-29);
\fill[diffSame] (19,-28) rectangle (20,-29);
\fill[diffDiff] (20,-28) rectangle (21,-29);
\fill[diffSame] (21,-28) rectangle (22,-29);
\fill[diffDiff] (22,-28) rectangle (24,-29);
\fill[diffSame] (0,-29) rectangle (17,-30);
\fill[diffDiff] (17,-29) rectangle (18,-30);
\fill[diffSame] (18,-29) rectangle (19,-30);
\fill[diffDiff] (19,-29) rectangle (20,-30);
\fill[diffSame] (20,-29) rectangle (21,-30);
\fill[diffDiff] (21,-29) rectangle (24,-30);
\fill[diffSame] (0,-30) rectangle (16,-31);
\fill[diffDiff] (16,-30) rectangle (17,-31);
\fill[diffSame] (17,-30) rectangle (18,-31);
\fill[diffDiff] (18,-30) rectangle (19,-31);
\fill[diffSame] (19,-30) rectangle (20,-31);
\fill[diffDiff] (20,-30) rectangle (24,-31);
\fill[diffSame] (0,-31) rectangle (15,-32);
\fill[diffDiff] (15,-31) rectangle (18,-32);
\fill[diffSame] (18,-31) rectangle (19,-32);
\fill[diffDiff] (19,-31) rectangle (24,-32);
\fill[diffSame] (0,-32) rectangle (14,-33);
\fill[diffDiff] (14,-32) rectangle (17,-33);
\fill[diffSame] (17,-32) rectangle (18,-33);
\fill[diffDiff] (18,-32) rectangle (24,-33);
\fill[diffSame] (0,-33) rectangle (13,-34);
\fill[diffDiff] (13,-33) rectangle (24,-34);
\fill[diffSame] (0,-34) rectangle (12,-35);
\fill[diffDiff] (12,-34) rectangle (18,-35);
\fill[diffSame] (18,-34) rectangle (19,-35);
\fill[diffDiff] (19,-34) rectangle (24,-35);
\fill[diffSame] (0,-35) rectangle (11,-36);
\fill[diffDiff] (11,-35) rectangle (24,-36);
\fill[diffSame] (0,-36) rectangle (10,-37);
\fill[diffDiff] (10,-36) rectangle (19,-37);
\fill[diffSame] (19,-36) rectangle (21,-37);
\fill[diffDiff] (21,-36) rectangle (24,-37);
\fill[diffSame] (0,-37) rectangle (9,-38);
\fill[diffDiff] (9,-37) rectangle (19,-38);
\fill[diffSame] (19,-37) rectangle (22,-38);
\fill[diffDiff] (22,-37) rectangle (24,-38);
\fill[diffSame] (0,-38) rectangle (8,-39);
\fill[diffDiff] (8,-38) rectangle (11,-39);
\fill[diffSame] (11,-38) rectangle (12,-39);
\fill[diffDiff] (12,-38) rectangle (19,-39);
\fill[diffSame] (19,-38) rectangle (22,-39);
\fill[diffDiff] (22,-38) rectangle (24,-39);
\fill[diffSame] (0,-39) rectangle (7,-40);
\fill[diffDiff] (7,-39) rectangle (19,-40);
\fill[diffSame] (19,-39) rectangle (23,-40);
\fill[diffDiff] (23,-39) rectangle (24,-40);
\fill[diffSame] (0,-40) rectangle (6,-41);
\fill[diffDiff] (6,-40) rectangle (17,-41);
\fill[diffSame] (17,-40) rectangle (18,-41);
\fill[diffDiff] (18,-40) rectangle (19,-41);
\fill[diffSame] (19,-40) rectangle (22,-41);
\fill[diffDiff] (22,-40) rectangle (23,-41);
\fill[diffSame] (23,-40) rectangle (24,-41);
\fill[diffSame] (0,-41) rectangle (5,-42);
\fill[diffDiff] (5,-41) rectangle (20,-42);
\fill[diffSame] (20,-41) rectangle (21,-42);
\fill[diffDiff] (21,-41) rectangle (22,-42);
\fill[diffSame] (22,-41) rectangle (23,-42);
\fill[diffDiff] (23,-41) rectangle (24,-42);
\fill[diffSame] (0,-42) rectangle (4,-43);
\fill[diffDiff] (4,-42) rectangle (7,-43);
\fill[diffSame] (7,-42) rectangle (8,-43);
\fill[diffDiff] (8,-42) rectangle (12,-43);
\fill[diffSame] (12,-42) rectangle (13,-43);
\fill[diffDiff] (13,-42) rectangle (17,-43);
\fill[diffSame] (17,-42) rectangle (18,-43);
\fill[diffDiff] (18,-42) rectangle (20,-43);
\fill[diffSame] (20,-42) rectangle (21,-43);
\fill[diffDiff] (21,-42) rectangle (24,-43);
\fill[diffSame] (0,-43) rectangle (3,-44);
\fill[diffDiff] (3,-43) rectangle (9,-44);
\fill[diffSame] (9,-43) rectangle (10,-44);
\fill[diffDiff] (10,-43) rectangle (17,-44);
\fill[diffSame] (17,-43) rectangle (18,-44);
\fill[diffDiff] (18,-43) rectangle (24,-44);
\fill[diffSame] (0,-44) rectangle (2,-45);
\fill[diffDiff] (2,-44) rectangle (6,-45);
\fill[diffSame] (6,-44) rectangle (7,-45);
\fill[diffDiff] (7,-44) rectangle (13,-45);
\fill[diffSame] (13,-44) rectangle (14,-45);
\fill[diffDiff] (14,-44) rectangle (15,-45);
\fill[diffSame] (15,-44) rectangle (16,-45);
\fill[diffDiff] (16,-44) rectangle (24,-45);
\fill[diffSame] (0,-45) rectangle (1,-46);
\fill[diffDiff] (1,-45) rectangle (2,-46);
\fill[diffSame] (2,-45) rectangle (3,-46);
\fill[diffDiff] (3,-45) rectangle (21,-46);
\fill[diffSame] (21,-45) rectangle (22,-46);
\fill[diffDiff] (22,-45) rectangle (24,-46);
\fill[diffDiff] (0,-46) rectangle (24,-47);

    \draw[white,line width=2.0pt] \staircase{46}{23} -- (24,-23);
    \draw[ink,densely dashed,line width=0.8pt] \staircase{46}{23} -- (24,-23);
    \draw[grid,line width=0.3pt] (0,0) rectangle (24,-47);
    \node[font=\small,anchor=south] at (12,0.6) {(c) (a) compared with (b)};
    \foreach \x/\l in {1/1,24/24}
      \node[font=\scriptsize,anchor=north] at (\x-0.5,-47.3) {$\l$};
  \end{scope}
  % legends
  \node[anchor=north,font=\scriptsize] at (52.5,-50.3) {in (a) and (b):\quad \swatch{white}\,$\Lq$ (quiescent)\quad
    \swatch{stG}\,$\Gq$ (general)\quad \swatch{stA}\,$\Aq$\quad \swatch{stB}\,$\Bq$\quad
    \swatch{stC}\,$\Cq$\quad \swatch{stF}\,$\Fq$ (firing)};
  \node[anchor=north,font=\scriptsize] at (52.5,-53.9) {in (c):\quad \swatch{white}\,quiescent in both,
    $t<i-1$\quad \swatch{diffSame}\,the same state in (a) and (b)\quad
    \swatch{diffDiff}\,different states};
\end{tikzpicture}

\caption{Mazoyer's six-state minimal-time solution. Left: the line of length
$n=24$; time runs downwards from $t=0$ to the firing time $t=2n-2=46$, and the
cells $i=1,\dots,24$ run from left to right. Right: the half-line, which is
started in the same way but has no right end, for the same times. The two
diagrams coincide above the dashed anti-diagonal $t+i=2n-2$
(\cref{lem:agree}); the dotted line marks the right end of the line of length
$24$. Below the anti-diagonal lies the reflected triangle $R_n$, in which the
line reacts to its right end; on the first anti-diagonal of $R_n$,
the \emph{return chain} $t+i=2n-1$, every cell of the line differs from the
corresponding cell of the half-line (\cref{lem:return}).}
\label{fig:line}
\end{figure}

\paragraph{Results}
Our main results concern minimal-time solutions with any number of states.
\begin{enumerate}[(1)]
\item \emph{The speed-$1/3$ barrier} (\cref{thm:third}). Behind the front of
  the half-line, each row parallel to the front starts at the left border and
  becomes periodic after a transient. Infinitely many of these rows become
  periodic only after crossing the line $i=t/3$, up to an arbitrarily small
  angle. The proof rests on a pumping lemma (\cref{lem:pumping}): two reflected
  triangles that receive the same input on the cone of their right end fire it
  at the same relative time, which is impossible for lines of different
  lengths. In the known constructions, the boundary of the periodic part is the
  signal that meets the echo at the first division point of the line; read in
  this way, the barrier says that this point is asymptotically never in the half
  of the line next to the general (\cref{rem:division}). In simulations, the
  half-line of the five-state rule $\delta_{12}$ of \cref{sec:five}, which
  synchronizes the lines of lengths $2,\dots,12$, meets the barrier exactly.
\item \emph{The left-border barrier} (\cref{thm:left,cor:regular}). The
  half-line is never periodic in time and space on a wedge at its left border;
  in particular it is never \emph{eventually regular}, that is, it never
  consists of finitely many periodic regions separated by signals of rational
  speeds. The reason is that the cells near the left border are fired by a
  finite automaton that reads the half-line along anti-diagonals, and a long
  periodic input would make it fire too early (\cref{lem:left}). The half-line
  of $\delta_{12}$, which meets the speed-$1/3$ barrier exactly, is excluded by
  this one (\cref{cor:germ12}).
\item \emph{Finitely many early rows} (\cref{cor:reduce}). The first rows of
  the reflected triangles of all lengths take only finitely many forms, which a
  certificate computed from the half-line determines. This turns the condition
  that no reflected triangle fires early in its first rows into finitely many
  checks (\cref{lem:G}).
\end{enumerate}
With these conditions and satisfiability solvers whose negative answers are
certified by independently checked proofs, we obtain the following results.
Four states are excluded by the lines of lengths $2,\dots,5$ together with
partial information on the lines of lengths $6,\dots,10$, and this is sharp in
several directions (\cref{thm:four,prop:sharp}). Five five-state rules
synchronize exactly the lines of lengths $2,\dots,14$
(\cref{prop:nearmiss}), so the lines of lengths $2,\dots,14$ alone cannot
exclude five states. On the cells of its half-line with $t+i\le78$, every
five-state minimal-time solution uses at least $15$ distinct neighbourhoods
besides the quiescent one (\cref{prop:germs}). Finally, a conditional five-state
result reported by Balzer in 1967 holds, even in a weaker form
(\cref{prop:balzer}). \Cref{tab:results} lists the results with the kind of
evidence behind each of them.

\begin{table}[t]
\centering\small
\caption{The results of this paper and the evidence behind them. A
``computation'' is a finite simulation or enumeration; ``certified'' means that
every unsatisfiability claim comes with a proof that an independent checker
verifies.}
\label{tab:results}
\begin{tabular}{@{}>{\raggedright\arraybackslash}p{0.56\textwidth}l>{\raggedright\arraybackslash}p{0.25\textwidth}@{}}
\toprule
Result & Where & Evidence\\
\midrule
Agreement with the half-line; return chain; determinacy & \ref{lem:agree}, \ref{lem:return}, \ref{lem:determinacy} & proof\\
Pumping lemma; speed-$1/3$ barrier & \ref{lem:pumping}, \ref{thm:third} & proof\\
Right band & \ref{lem:rightband} & proof\\
Left-border lemma and barrier; no eventually regular half-line & \ref{lem:left}, \ref{thm:left}, \ref{cor:regular} & proof\\
Finitely many early rows; the formulas $\G_k(T,N_1,C)$ & \ref{cor:reduce}, \ref{lem:G} & proof\\
Four states excluded by the lines of lengths $2,\dots,5$ and parts of those of lengths $6,\dots,10$ & \ref{thm:four} & certified computation\\
Sharpness of \cref{thm:four} & \ref{prop:sharp} & computation\\
Five rules for exactly the lengths $2,\dots,14$ & \ref{prop:nearmiss} & computation and proof\\
At least $15$ neighbourhoods on the half-line of a five-state solution & \ref{prop:germs} & certified computation and proof\\
The half-line of $\delta_{12}$ and $\delta_{13}$ cannot be completed & \ref{cor:germ12} & computation and proof\\
Balzer's conditional five-state result & \ref{prop:balzer} & certified computation\\
\bottomrule
\end{tabular}
\end{table}

\paragraph{Related work}
Indistinguishability arguments, in which a solution cannot tell two situations
apart and therefore fires too early, show that there is no minimal-time
solution with one-way communication~\cite[Prop.~2]{Mazoyer96}. Through
repetitions forced by the pigeonhole principle, they also exclude minimal-time
solutions for variants with sub-generals~\cite{YamashitaEtAl14} and for some
geometric configurations~\cite{Kobayashi19}. Our pumping lemma and our
left-border lemma use the same mechanism to constrain solutions of the original
problem. That the behaviour behind the front of a cellular automaton is
eventually periodic is classical~\cite{MazoyerTerrier99}; our contribution is a
lower bound on where this periodicity can start. Lower bounds on the number of
states are known for restricted classes and for variants such as the
ring~\cite{BerthiaumeEtAl04}; see the survey~\cite{Correa17}. Satisfiability
solvers have been applied to cellular automata before~\cite{DAntonioDelzanno04},
and clausal proofs checked by programs independent of the
solvers~\cite{WetzlerHeuleHunt14,CruzFilipeEtAl17} have certified large
computer-assisted results in combinatorics~\cite{HeuleKullmannMarek16,Heule18}.
We are not aware
of statements equivalent to
\cref{thm:third,thm:left,cor:regular,cor:reduce} in the literature.

\paragraph{Organization}
\Cref{sec:lines} introduces lines, the half-line and their agreement.
\Cref{sec:right} studies the right end of the reflected triangles and proves
the speed-$1/3$ barrier, \cref{sec:left} studies their left end and proves the
left-border barrier, and \cref{sec:early} reduces their first rows to finitely
many cases. \Cref{sec:applications} applies these results to four and five
states, and \cref{sec:discussion} discusses open questions. A list of symbols
is given at the end. The programs, rule tables and logs behind every
computational statement, with scripts that rerun them, are in the
supplementary material, as described in the data availability statement at the end.
